\chapter{Tests et analyse des resultats}

\section{Protocole experimental}
L'evaluation est structuree en quatre blocs, alignes sur la chaine de
traitement complete:
\begin{itemize}
\item tests de compensation de mouvement (stabilite du fond apres warping);
\item tests de segmentation (qualite du masque foreground);
\item tests de clustering (coherence des objets detectes);
\item tests d'integration navigation (prise en compte des obstacles dynamiques
dans Nav2).
\end{itemize}

Les scenarios incluent des phases statiques, des rotations du robot, puis des
deplacements avec un ou plusieurs objets en mouvement. Les indicateurs suivis
restent volontairement simples:
\begin{itemize}
\item densite visuelle des evenements compenses sur le fond;
\item taux de bruit residuel dans le masque;
\item stabilite des contours d'objets sur des sequences consecutives;
\item impact sur la trajectoire Nav2 (detours, arrets de securite, absence de
collision).
\end{itemize}

\section{Resultats}
Les essais fonctionnels confirment le comportement attendu de la pipeline:
\begin{itemize}
\item la compensation IMU reduit la dispersion des evenements de fond et rend
les zones dynamiques plus discriminantes;
\item la segmentation par seuil dynamique produit un masque exploitable meme
quand la scene contient des changements rapides;
\item le clustering permet de convertir le masque en objets utilisables pour la
navigation.
\end{itemize}

Sur le volet navigation, l'ajout de la source
\texttt{/dynamic\_obstacles} dans les costmaps locale et globale permet de
refleter les obstacles detectes dans la carte de cout et d'influencer le plan
local en temps reel.

\section{Analyse}
L'analyse met en evidence trois points:
\begin{itemize}
\item le couplage IMU + evenement est essentiel: sans compensation, la
segmentation degrade rapidement des que l'ego-mouvement augmente;
\item le seuil dynamique $\lambda$ est un parametre sensible qui conditionne le
compromis entre bruit et rappel d'objets;
\item l'integration Nav2 via les costmaps est robuste tant que le topic
\texttt{/dynamic\_obstacles} est publie de facon stable.
\end{itemize}

Les limites actuelles concernent surtout les cas de faible texture, les objets
tres proches les uns des autres, et les scenes a fort bruit evenementiel.

\section{Discussion}
La solution retenue est coherente avec l'objectif du projet: obtenir une
perception d'obstacles dynamiques compatible avec des contraintes temps reel sur
robot mobile.

Les pistes d'amelioration prioritaires sont:
\begin{itemize}
\item calibration plus fine des parametres de compensation et de segmentation
selon les conditions lumineuses;
\item enrichissement du clustering (metriques mouvement plus explicites);
\item etape supplementaire d'association multi-frames pour stabiliser les objets
et estimer leur vitesse avant projection dans Nav2.
\end{itemize}
